\begin{figure}[htbp]
\centering
\begin{tikzpicture}[>=stealth,x=.9cm]
\draw[->,black!50] (-3.3,0)--(3.4,0) node[right] {$x$};
\draw[arkivblaa] (-2,-.08)--(-2,.08) (0,-.08)--(0,.08) (2,-.08)--(2,.08);
\node[below] at (-2,0) {$x_0-r$};\node[below] at (0,0) {$x_0$};
\node[below] at (2,0) {$x_0+r$};
\draw[<->,arkivblaa] (-2,.45)--(0,.45) node[midway,above] {$r$};
\draw[<->,arkivblaa] (0,.45)--(2,.45) node[midway,above] {$r$};
\node at (0,1.25) {$|x-x_0|=r$};
\draw[thick,draw=arkivblaa,fill=white] (-2,0) circle(2pt);
\draw[thick,draw=arkivblaa,fill=white] (2,0) circle(2pt);
\end{tikzpicture}
\caption{Absolutværdi som afstand: De to punkter har samme afstand $r>0$ til $x_0$. Betingelsen $|x-x_0|<r$ vælger punkterne imellem dem og udelader endepunkterne. Talaksen er skematisk og dimensionsløs.}
\label{fig:b1-afstand}
\end{figure}
